% Sorgente LaTeX canonico, importato e revisionato dal precedente sorgente Markdown.
\chapter{Funzioni e geometria analitica}
\label{chap:3}

\begin{obiettivocapitolo}{determinare dominio, inversa, trasformazioni e oggetti geometrici}
\prioritaalta. Il grafico e il dominio anticipano spesso il comportamento richiesto in limiti, derivate e integrali.
\end{obiettivocapitolo}
\begin{proceduraesame}{dominio di una funzione composta}
\begin{enumerate}
\item Elenca tutte le condizioni: denominatori $\ne0$, radicandi pari $\ge0$, argomenti logaritmici $>0$, argomenti di inverse trigonometriche negli intervalli ammessi.
\item Risolvi ciascuna condizione e interseca gli insiemi ottenuti.
\item Per $g\circ f$, richiedi contemporaneamente $x\in D_f$ e $f(x)\in D_g$.
\end{enumerate}
\end{proceduraesame}
\begin{proceduraesame}{funzione inversa}
\begin{enumerate}
\item Verifica o imponi l'iniettività, eventualmente restringendo il dominio a un intervallo monotono.
\item Scrivi $y=f(x)$, risolvi rispetto a $x$, scegli il ramo compatibile con il dominio ristretto.
\item Scambia il nome delle variabili e dichiara dominio e immagine dell'inversa.
\item Controlla $f^{-1}(f(x))=x$ nel dominio e $f(f^{-1}(y))=y$ nell'immagine.
\end{enumerate}
\end{proceduraesame}
\begin{sceltametodo}{trasformazioni del grafico}
Applica prima le trasformazioni interne all'argomento e poi quelle esterne al valore. In $y=a f(b(x-h))+k$, gestisci $x-h$, poi la scala/riflessione orizzontale $b$, quindi scala verticale $a$ e traslazione $k$.
\end{sceltametodo}
\begin{controllofinale}{geometria analitica}
Verifica che il punto sostituito soddisfi l'equazione; controlla pendenza o vettore normale; nelle distanze il risultato deve essere non negativo.
\end{controllofinale}

\section{Funzioni elementari di riferimento}\label{funzioni-elementari-di-riferimento}

Per \(n,m\in\mathbb N_+\), salvo diversa indicazione, le funzioni di uso frequente sono:

\begin{longtable}[]{@{}
  >{\raggedright\arraybackslash}p{(\columnwidth - 6\tabcolsep) * \real{0.2500}}
  >{\raggedright\arraybackslash}p{(\columnwidth - 6\tabcolsep) * \real{0.2500}}
  >{\raggedright\arraybackslash}p{(\columnwidth - 6\tabcolsep) * \real{0.2500}}
  >{\raggedright\arraybackslash}p{(\columnwidth - 6\tabcolsep) * \real{0.2500}}@{}}
\toprule\noalign{}
\begin{minipage}[b]{\linewidth}\raggedright
Funzione
\end{minipage} & \begin{minipage}[b]{\linewidth}\raggedright
Dominio
\end{minipage} & \begin{minipage}[b]{\linewidth}\raggedright
Immagine
\end{minipage} & \begin{minipage}[b]{\linewidth}\raggedright
Proprietà essenziale
\end{minipage} \\
\midrule\noalign{}
\endhead
\bottomrule\noalign{}
\endlastfoot
\(x^n\), \(n\) pari & \(\mathbb R\) & \([0,+\infty)\) & pari, decrescente su \((-\infty,0]\), crescente su \([0,+\infty)\) \\
\(x^n\), \(n\) dispari & \(\mathbb R\) & \(\mathbb R\) & dispari, strettamente crescente \\
\(\lvert x\rvert\) & \(\mathbb R\) & \([0,+\infty)\) & pari, decrescente su \((-\infty,0]\), crescente su \([0,+\infty)\) \\
\(1/x\) & \(\mathbb R\setminus\{0\}\) & \(\mathbb R\setminus\{0\}\) & dispari, strettamente decrescente su ciascuna componente del dominio \\
\(\sqrt[2m]{x}\) & \([0,+\infty)\) & \([0,+\infty)\) & strettamente crescente \\
\(\sqrt[2m+1]{x}\) & \(\mathbb R\) & \(\mathbb R\) & strettamente crescente \\
\(a^x\), \(a>1\) & \(\mathbb R\) & \((0,+\infty)\) & strettamente crescente \\
\(a^x\), \(0<a<1\) & \(\mathbb R\) & \((0,+\infty)\) & strettamente decrescente \\
\(\ln x\) & \((0,+\infty)\) & \(\mathbb R\) & strettamente crescente \\
\end{longtable}

Parte intera, parte frazionaria e segno:

\[
\lfloor x\rfloor\le x<\lfloor x\rfloor+1,
\qquad
\{x\}=x-\lfloor x\rfloor\in[0,1),
\]

\[
\operatorname{sgn}x=
\begin{cases}
-1,&x<0,\\
0,&x=0,\\
1,&x>0.
\end{cases}
\]

\section{Dominio, immagine, preimmagine, composizione e inversa}\label{dominio-immagine-preimmagine-composizione-e-inversa}

Nella scrittura \(f:A\to B\), \(A\) è il dominio e \(B\) il codominio. L'immagine \(f(A)\) può essere un sottoinsieme proprio del codominio; la suriettività richiede \(f(A)=B\).

Il grafico e la preimmagine di un insieme \(Y\subseteq B\) sono

\[
G_f=\{(x,f(x)):x\in A\},
\qquad
f^{-1}(Y)=\{x\in A:f(x)\in Y\}.
\]

La notazione \(f^{-1}(Y)\) indica una \textbf{preimmagine} ed è definita anche quando \(f\) non è invertibile. La funzione inversa \(f^{-1}:B\to A\) esiste invece solo quando \(f\) è biiettiva.

\[
f:A\to B,
\qquad
f(A)=\{f(x):x\in A\}\subseteq B.
\]

\[
f\text{ iniettiva}
\Longleftrightarrow
f(x_1)=f(x_2)\Rightarrow x_1=x_2,
\]

\[
f\text{ suriettiva su }B
\Longleftrightarrow
\forall y\in B,\ \exists x\in A:\ f(x)=y,
\]

\[
f\text{ biiettiva}
\Longleftrightarrow
f\text{ iniettiva e suriettiva}.
\]

Se \(f:A\to B\) è biiettiva, per \(x\in A\) e \(y\in B\):

\[
f^{-1}(y)=x
\Longleftrightarrow
f(x)=y.
\]

\[
f^{-1}\circ f=\operatorname{id}_A,
\qquad
f\circ f^{-1}=\operatorname{id}_B.
\]

\[
(g\circ f)(x)=g(f(x)),
\qquad
D_{g\circ f}=\{x\in D_f:f(x)\in D_g\}.
\]

\[
(g\circ f)^{-1}=f^{-1}\circ g^{-1},
\]

quando \(f\) e \(g\) sono invertibili.

Coppie diretta/inversa frequenti:

\begin{longtable}[]{@{}
  >{\raggedright\arraybackslash}p{(\columnwidth - 4\tabcolsep) * \real{0.3333}}
  >{\raggedright\arraybackslash}p{(\columnwidth - 4\tabcolsep) * \real{0.3333}}
  >{\raggedright\arraybackslash}p{(\columnwidth - 4\tabcolsep) * \real{0.3333}}@{}}
\toprule\noalign{}
\begin{minipage}[b]{\linewidth}\raggedright
Funzione diretta
\end{minipage} & \begin{minipage}[b]{\linewidth}\raggedright
Restrizione che la rende biiettiva
\end{minipage} & \begin{minipage}[b]{\linewidth}\raggedright
Funzione inversa
\end{minipage} \\
\midrule\noalign{}
\endhead
\bottomrule\noalign{}
\endlastfoot
\(y=x^n\), \(n\) pari & \(x\in[0,+\infty)\) & \(x=\sqrt[n]{y}\), \(y\ge0\) \\
\(y=x^n\), \(n\) pari & \(x\in(-\infty,0]\) & \(x=-\sqrt[n]{y}\), \(y\ge0\) \\
\(y=x^n\), \(n\) dispari & \(x\in\mathbb R\) & \(x=\sqrt[n]{y}\) \\
\(y=a^x\), \(a>0\), \(a\ne1\) & \(x\in\mathbb R\) & \(x=\log_a y\), \(y>0\) \\
\(y=\sin x\) & \(x\in[-\pi/2,\pi/2]\) & \(x=\arcsin y\), \(\lvert y\rvert\le1\) \\
\(y=\cos x\) & \(x\in[0,\pi]\) & \(x=\arccos y\), \(\lvert y\rvert\le1\) \\
\(y=\tan x\) & \(x\in(-\pi/2,\pi/2)\) & \(x=\arctan y\) \\
\end{longtable}

Per \(D_f=-D_f\):

\[
f(-x)=f(x)\quad\text{pari},
\qquad
f(-x)=-f(x)\quad\text{dispari}.
\]

\[
D_f+T=D_f,
\qquad
f(x+T)=f(x),
\qquad T>0.
\]

Qui \(D_f+T=\{x+T:x\in D_f\}\). Un \textbf{periodo fondamentale} è il minimo periodo positivo, quando tale minimo esiste.

Limitatezza su \(E\subseteq D_f\):

\[
\begin{aligned}
f\text{ limitata superiormente su }E
&\Longleftrightarrow \exists M\in\mathbb R:\ f(x)\le M\quad\forall x\in E,\\
f\text{ limitata inferiormente su }E
&\Longleftrightarrow \exists m\in\mathbb R:\ m\le f(x)\quad\forall x\in E,\\
f\text{ limitata su }E
&\Longleftrightarrow \exists C>0:\ |f(x)|\le C\quad\forall x\in E.
\end{aligned}
\]

Monotonia su un intervallo \(I\):

\[
\begin{aligned}
x_1<x_2&\Longrightarrow f(x_1)\le f(x_2)
&&\text{non decrescente},\\
x_1<x_2&\Longrightarrow f(x_1)<f(x_2)
&&\text{strettamente crescente},\\
x_1<x_2&\Longrightarrow f(x_1)\ge f(x_2)
&&\text{non crescente},\\
x_1<x_2&\Longrightarrow f(x_1)>f(x_2)
&&\text{strettamente decrescente}.
\end{aligned}
\]

Estremi assoluti su \(E\):

\[
\begin{aligned}
x_M\in E\text{ è punto di massimo assoluto}
&\Longleftrightarrow f(x)\le f(x_M)\quad\forall x\in E,\\
x_m\in E\text{ è punto di minimo assoluto}
&\Longleftrightarrow f(x_m)\le f(x)\quad\forall x\in E.
\end{aligned}
\]

Per un estremo locale le stesse disuguaglianze sono richieste soltanto in un intorno relativo del punto considerato.

Approfondimento: \href{https://ingegnerismo.it/matematica/funzione-matematica/}{funzione matematica}, \href{https://ingegnerismo.it/matematica/funzione-inversa/}{funzione inversa}, \href{https://ingegnerismo.it/matematica/funzione-monotona/}{funzione monotona}.

\section{Trasformazioni dei grafici}\label{trasformazioni-dei-grafici}

\[
\begin{array}{c|c}
\text{funzione}&\text{trasformazione di }y=f(x)\\
\hline
f(x)+k&\text{traslazione verticale di }k\\
f(x-h)&\text{traslazione orizzontale di }h\\
af(x)&\text{scala verticale }|a|;\ \text{riflessione rispetto all'asse }x\text{ se }a<0\\
f(bx)&\text{scala orizzontale }1/|b|\ (b\ne0);\ \text{riflessione rispetto all'asse }y\text{ se }b<0\\
-f(x)&\text{riflessione rispetto all'asse }x\\
f(-x)&\text{riflessione rispetto all'asse }y\\
|f(x)|&\text{rami negativi riflessi sopra l'asse }x\\
f(|x|)&\text{ramo }x\ge0\text{ copiato per }x\le0
\end{array}
\]

Per \(b=0\), \(f(bx)=f(0)\) soltanto se \(0\in D_f\).

Se \(f\) è biiettiva, il grafico dell'inversa si ottiene riflettendo il grafico di \(f\) rispetto alla bisettrice \(y=x\):

\[
G_{f^{-1}}=\text{riflessione di }G_f\text{ rispetto a }y=x.
\]

Approfondimento: \href{https://ingegnerismo.it/matematica/trasformazioni-di-grafici/}{trasformazioni di grafici}.

\section{Geometria analitica nel piano}\label{geometria-analitica-nel-piano}

Le coordinate sono espresse in un sistema cartesiano ortonormale. Per \(u=(u_1,u_2)\) e \(v=(v_1,v_2)\):

\[
u\cdot v=u_1v_1+u_2v_2,
\qquad
\lVert u\rVert=\sqrt{u_1^2+u_2^2},
\]

\[
u\perp v\Longleftrightarrow u\cdot v=0,
\qquad
\cos\theta=\frac{u\cdot v}{\lVert u\rVert\,\lVert v\rVert}
\quad(u,v\ne0),
\]

con \(\theta\in[0,\pi]\) angolo non orientato fra i due vettori.

Retta passante per \(P_0=(x_0,y_0)\) con vettore direttore \(v=(v_1,v_2)\ne0\):

\[
(x,y)=(x_0,y_0)+t(v_1,v_2),
\qquad t\in\mathbb R.
\]

Per \(A(x_1,y_1)\) e \(B(x_2,y_2)\), \(AB\) indica la distanza e \(M\) il punto medio del segmento:

\[
AB=\sqrt{(x_2-x_1)^2+(y_2-y_1)^2},
\]

\[
M=\left(\frac{x_1+x_2}{2},\frac{y_1+y_2}{2}\right).
\]

\[
m=\frac{y_2-y_1}{x_2-x_1},
\qquad x_1\ne x_2,
\]

\[
y-y_1=m(x-x_1),
\qquad
x=x_1\quad(x_1=x_2).
\]

Forma generale, forma esplicita e forma degli intercetti:

\[
ax+by+c=0,
\qquad (a,b)\ne(0,0),
\]

\[
y=mx+q,
\qquad
m=-\frac ab,
\qquad
q=-\frac cb,
\qquad b\ne0,
\]

\[
\frac xp+\frac yq=1,
\qquad pq\ne0,
\]

in cui \(p\) e \(q\) sono rispettivamente le intercette sugli assi \(x\) e \(y\). Dalla forma generale:

\[
p=-\frac ca,
\qquad
q=-\frac cb,
\qquad
abc\ne0.
\]

Se \(c=0\), la retta passa per l'origine e la forma degli intercetti con \(pq\ne0\) non è disponibile. Se \(a=0\) oppure \(b=0\), la retta è parallela a uno degli assi e manca la corrispondente intercetta finita.

\[
r_1\parallel r_2
\Longleftrightarrow
a_1b_2-a_2b_1=0,
\]

\[
r_1\perp r_2
\Longleftrightarrow
a_1a_2+b_1b_2=0.
\]

La prima condizione confronta le direzioni e comprende anche il caso di rette coincidenti; per richiedere rette parallele distinte occorre escludere la coincidenza.

Se le rette non sono verticali, con coefficienti angolari \(m_1,m_2\):

\[
r_1\parallel r_2\Longleftrightarrow m_1=m_2,
\qquad
r_1\perp r_2\Longleftrightarrow m_1m_2=-1.
\]

Per \(P=(x_0,y_0)\) e \(r:ax+by+c=0\):

\[
d(P,r)=\frac{|ax_0+by_0+c|}{\sqrt{a^2+b^2}}.
\]

Circonferenza:

\[
(x-h)^2+(y-k)^2=r^2,
\qquad r\ge0,
\]

\[
x^2+y^2+Dx+Ey+F=0.
\]

Passaggio fra le due forme:

\[
D=-2h,
\qquad
E=-2k,
\qquad
F=h^2+k^2-r^2,
\]

\[
C=\left(-\frac D2,-\frac E2\right),
\qquad
r^2=\frac{D^2+E^2}{4}-F.
\]

\begin{longtable}[]{@{}ll@{}}
\toprule\noalign{}
Condizione & Luogo reale \\
\midrule\noalign{}
\endhead
\bottomrule\noalign{}
\endlastfoot
\(r^2>0\) & circonferenza di raggio \(r=\sqrt{r^2}\) \\
\(r^2=0\) & il solo punto \(C\) \\
\(r^2<0\) & insieme vuoto in \(\mathbb R^2\) \\
\end{longtable}

Parabola con asse verticale:

\[
y=ax^2+bx+c,
\qquad a\ne0,
\]

\[
x_V=-\frac b{2a},
\qquad
y_V=-\frac\Delta{4a},
\qquad
\text{asse}:x=x_V,
\]

\[
y=a(x-x_V)^2+y_V.
\]

Dalla forma del vertice alla forma esplicita:

\[
b=-2ax_V,
\qquad
c=ax_V^2+y_V.
\]

La parabola è rivolta verso l'alto per \(a>0\) e verso il basso per \(a<0\).

Coniche canoniche non ruotate. In tutte le formule \((h,k)\) è il centro; \(a\) indica il semiasse trasverso o maggiore e \(b\) l'altro semiasse.

Ellisse con asse maggiore orizzontale:

\[
\frac{(x-h)^2}{a^2}+\frac{(y-k)^2}{b^2}=1,
\qquad a,b>0,
\]

\[
c^2=a^2-b^2,
\qquad a\ge b,
\qquad
F_{1,2}=(h\pm c,k).
\]

Iperbole con asse trasverso orizzontale:

\[
\frac{(x-h)^2}{a^2}-\frac{(y-k)^2}{b^2}=1,
\qquad a,b>0,
\]

\[
c^2=a^2+b^2,
\qquad
F_{1,2}=(h\pm c,k),
\qquad
y-k=\pm\frac ba(x-h).
\]

Ellisse con asse maggiore verticale:

\[
\frac{(x-h)^2}{b^2}+\frac{(y-k)^2}{a^2}=1,
\qquad
c^2=a^2-b^2,
\qquad
F_{1,2}=(h,k\pm c),
\qquad a\ge b>0.
\]

Iperbole con asse trasverso verticale:

\[
\frac{(y-k)^2}{a^2}-\frac{(x-h)^2}{b^2}=1,
\qquad
c^2=a^2+b^2,
\qquad
F_{1,2}=(h,k\pm c),
\]

\[
y-k=\pm\frac ab(x-h),
\qquad a,b>0.
\]

Approfondimento: \href{https://ingegnerismo.it/matematica/retta/}{retta}, \href{https://ingegnerismo.it/matematica/circonferenza/}{circonferenza}, \href{https://ingegnerismo.it/matematica/parabola/}{parabola}, \href{https://ingegnerismo.it/matematica/ellisse/}{ellisse}, \href{https://ingegnerismo.it/matematica/iperbole/}{iperbole}.
